\section{Introduction}
\label{sec:intro}

% ---------------------------------------------------------------------------
% §1 Background — multi-branch agentic execution
% ---------------------------------------------------------------------------
Agentic LLM applications interleave generation with branching.
A coding agent proposes several edits or tool calls for one repository context and executes one~\cite{react,sweagent,openhands}; the same pattern appears when an inference-time controller expands a Tree-of-Thoughts frontier or a self-consistency pool~\cite{tot,selfconsistency,cot}.
At each step a committed trunk fans out into $k$ residuals, and one becomes the next spine while the rest are discarded~\cite{gsm8k,math,esc}.

% ---------------------------------------------------------------------------
% §2 Why important — compute / communication cost under PD disaggregation
% ---------------------------------------------------------------------------
Branching improves search coverage, but it also multiplies serving cost.
Under prefill--decode (PD) disaggregation, prompt processing and token generation run on separate instances, and the KV pages produced by prefill must cross a connector before decode can begin~\cite{distserve,splitwise}.
Every admitted residual may therefore pay residual prefill, KV transfer, and decode before it is discarded, and that cost accumulates as the agent repeats the fan-out.

% ---------------------------------------------------------------------------
% §3 Existing gap — path pruning vs.\ KV optimization
% ---------------------------------------------------------------------------
Prior work improves efficiency mainly along two axes, and neither removes this cost.
Decode-time path pruners act after generation has started~\cite{esc,specrej,dpts,stoploss}, so residual prefill and the connector transfer have often already been paid.
Prefix caching shares the trunk when its pages are resident~\cite{vllm,sglang}, but neither rejects a residual nor controls which pages cross the connector.
Even a fully shared trunk still pays for each residual suffix that is later discarded.
Pruning therefore remains \emph{stage-local}: each stage can discard only work that has already reached it.
The goal is a disaggregated design across prefill and decode, a hardware-algorithm co-design for efficient agentic LLM serving.

\begin{figure*}[t]
\centering
\resizebox{0.78\textwidth}{!}{% PD-Prune runtime in DistServe Fig.6 visual language.
% Flow: Requests → Controller → Prefill/Decoding Pruner → Instance.
\begin{tikzpicture}[
  x=1cm, y=1cm,
  font=\sffamily,
  line cap=round, line join=round,
]
\definecolor{dsctrl}{RGB}{233,150,170}
\definecolor{dsmodel}{RGB}{245,215,95}
\definecolor{dsrun}{RGB}{170,220,230}
\definecolor{dsgpu}{RGB}{95,185,115}
\definecolor{dsinst}{RGB}{235,235,235}
\definecolor{dsedge}{RGB}{55,55,55}
\definecolor{dspre}{RGB}{255,200,140}
\definecolor{dsdec}{RGB}{180,210,255}

\path (0,0.75) rectangle (16.2,9.0);

% ---- Requests ----
\foreach \x in {1.35, 2.15, 2.95, 3.75, 4.55} {
  \draw[-{Latex[length=2.4mm,width=1.9mm]}, line width=1.0pt, draw=dsedge]
    (\x,8.75) -- (\x,8.15);
}
\node[font=\footnotesize\sffamily] at (2.95,8.92) {Requests};

% ---- Controller ----
\draw[draw=dsedge, fill=dsctrl, line width=0.8pt]
  (0.70,7.35) rectangle (15.50,8.05);
\node[font=\large\sffamily\bfseries] at (8.10,7.82) {Controller};
\node[font=\scriptsize\sffamily] at (8.10,7.50) {decide before any kernel: residency $m_i^{\Pin},m_i^{\Din}$ and score $s_i$ $\rightarrow$ $a_i\in\{\Hit,\Admit,\Probe,\Reject\}$};

% ---- Controller → Pruners ----
\draw[-{Latex[length=2.4mm,width=1.9mm]}, line width=1.05pt, draw=dsedge]
  (3.15,7.35) -- (3.15,6.85);
\draw[-{Latex[length=2.4mm,width=1.9mm]}, line width=1.05pt, draw=dsedge]
  (13.05,7.35) -- (13.05,6.85);

% ---- Prefill Pruner (left small box) ----
\draw[draw=dsedge, fill=dspre, line width=0.75pt]
  (1.15,6.05) rectangle (5.15,6.75);
\node[font=\small\sffamily\bfseries] at (3.15,6.40) {Prefill Pruner};

% ---- Decoding Pruner (right small box) ----
\draw[draw=dsedge, fill=dsdec, line width=0.75pt]
  (11.05,6.05) rectangle (15.05,6.75);
\node[font=\small\sffamily\bfseries] at (13.05,6.40) {Decoding Pruner};

% ---- Pruners → Instances ----
\draw[-{Latex[length=2.4mm,width=1.9mm]}, line width=1.05pt, draw=dsedge]
  (3.15,6.05) -- (3.15,5.55);
\draw[-{Latex[length=2.4mm,width=1.9mm]}, line width=1.05pt, draw=dsedge]
  (13.05,6.05) -- (13.05,5.55);

% ===================== Prefill Instance (stacked) =====================
\draw[draw=dsedge, fill=dsinst, line width=0.7pt]
  (1.15,0.85) rectangle (6.15,5.05);
\draw[draw=dsedge, fill=dsinst, line width=0.7pt]
  (0.90,1.10) rectangle (5.90,5.30);
\draw[draw=dsedge, fill=dsinst, line width=0.85pt]
  (0.65,1.35) rectangle (5.65,5.55);

\node[font=\normalsize\sffamily\bfseries] at (3.15,5.25) {Prefill Instance};

\draw[draw=dsedge, fill=dsmodel, line width=0.7pt]
  (1.05,4.35) rectangle (5.25,4.95);
\node[font=\small\sffamily] at (3.15,4.65) {LLM Model};

\draw[draw=dsedge, fill=dsrun, line width=0.7pt]
  (1.05,1.65) rectangle (5.25,4.15);
\node[font=\small\sffamily] at (3.15,3.85) {Parallel Runtime};

\foreach \gx/\gy in {1.35/2.85, 3.25/2.85, 1.35/1.90, 3.25/1.90} {
  \draw[draw=dsedge, fill=dsgpu, line width=0.6pt]
    (\gx,\gy) rectangle ({\gx+1.50},{\gy+0.72});
  \node[font=\small\sffamily] at ({\gx+0.75},{\gy+0.36}) {GPU};
}

% ===================== Decoding Instance (stacked) =====================
\draw[draw=dsedge, fill=dsinst, line width=0.7pt]
  (11.05,0.85) rectangle (16.05,5.05);
\draw[draw=dsedge, fill=dsinst, line width=0.7pt]
  (10.80,1.10) rectangle (15.80,5.30);
\draw[draw=dsedge, fill=dsinst, line width=0.85pt]
  (10.55,1.35) rectangle (15.55,5.55);

\node[font=\normalsize\sffamily\bfseries] at (13.05,5.25) {Decoding Instance};

\draw[draw=dsedge, fill=dsmodel, line width=0.7pt]
  (10.95,4.35) rectangle (15.15,4.95);
\node[font=\small\sffamily] at (13.05,4.65) {LLM Model};

\draw[draw=dsedge, fill=dsrun, line width=0.7pt]
  (10.95,1.65) rectangle (15.15,4.15);
\node[font=\small\sffamily] at (13.05,3.85) {Parallel Runtime};

\foreach \gx/\gy in {11.25/2.85, 13.15/2.85, 11.25/1.90, 13.15/1.90} {
  \draw[draw=dsedge, fill=dsgpu, line width=0.6pt]
    (\gx,\gy) rectangle ({\gx+1.50},{\gy+0.72});
  \node[font=\small\sffamily] at ({\gx+0.75},{\gy+0.36}) {GPU};
}

% ===================== Thick hollow KV arrow =====================
\fill[white]
  (5.95,2.85) -- (8.35,2.85) -- (8.35,2.45) -- (9.85,3.20) --
  (8.35,3.95) -- (8.35,3.55) -- (5.95,3.55) -- cycle;
\draw[draw=dsedge, line width=1.25pt]
  (5.95,2.85) -- (8.35,2.85) -- (8.35,2.45) -- (9.85,3.20) --
  (8.35,3.95) -- (8.35,3.55) -- (5.95,3.55) -- cycle;

\node[font=\footnotesize\sffamily\bfseries] at (7.40,4.20) {KV Cache Transfer};
\node[font=\scriptsize\sffamily, align=center] at (7.40,2.05)
  {decoder-missing pages\\[-1pt]of admitted residuals only};

\end{tikzpicture}
}
\caption{\textbf{\sys{} runtime.}
The controller reads the residency of each residual on the prefill instance ($m_i^{\Pin}$) and the decode instance ($m_i^{\Din}$) together with a lightweight model-free heuristic score $s_i$ (the prefill-side score), and assigns an action $a_i\in\{\Hit,\Admit,\Probe,\Reject\}$ before any kernel runs.
The prefill instance computes only the non-resident suffix of admitted and probed residuals; the connector carries only the pages the decode instance lacks.
The decode-side pruner extends a probed residual beyond $t_0$ tokens only if its generated prefix clears $\tau_{\mathrm{dec}}$.
After commit, only the spine $x_\star$ is retained.}
\label{fig:system}
\end{figure*}

% ---------------------------------------------------------------------------
% §4 Challenges — First/Second/Third draft
% ---------------------------------------------------------------------------
% \paragraph{Challenges.}
% Extending pruning beyond the decode stage in a disaggregated serving system introduces three intertwined challenges.
% \textit{First}, identifying unnecessary work before execution is difficult because of limited information and distributed KV-cache residency.
% Specifically, the value of a candidate remains uncertain before further generation, while its KV cache may reside on the prefill instance but be absent from the decode instance, so the missing pages still require cross-instance transfer.
% This calls for lightweight decisions based on available candidate information and KV residency across both instances.
% \textit{Second}, early pruning introduces a trade-off between efficiency and reasoning quality.
% Aggressive pruning reduces computation and communication but risks discarding candidates that could eventually produce correct answers, whereas conservative pruning preserves such candidates at the cost of additional execution.
% This trade-off motivates selective exploration of uncertain candidates to improve pruning reliability.
% \textit{Third}, dynamic pruning complicates KV-cache management, because branches share KV state but have different lifetimes.
% Discarding a branch requires reclaiming its private KV without affecting shared state, while the selected path needs to be retained for subsequent agent steps.
% This highlights the need for coordinated KV sharing, reclamation, and reuse.

% ---------------------------------------------------------------------------
% §4 Challenges — pre-execution / efficiency / KV-lifetime draft
% ---------------------------------------------------------------------------
% Cutting a residual before its prefill and its connector transfer raises three challenges.
%
% \paragraph{Challenge 1: Low-overhead pre-execution decision.}
% Before any kernel runs, the controller has only the residual text and the KV residency on the prefill and decode instances.
% A page already resident on the prefill instance may still be missing on the decode instance, in which case the connector must transfer it.
% From those signals the controller decides whether the residual is skipped, prefilled, transferred, or decoded, and that decision has to stay off the GPU critical path.
%
% \paragraph{Challenge 2: Efficiency versus reasoning quality.}
% Rejecting a residual before prefill removes its prefill, its KV transfer, and its decode, but a weak pre-execution score can still belong to a residual that later reaches a correct answer.
% Aggressive rejection therefore cuts cost and can lose committed-answer accuracy; retaining every residual preserves that accuracy and pays the full fan-out.
% Residuals in between need a short subsequent generation before the system spends a full decode budget on them.
%
% \paragraph{Challenge 3: Shared KV with unequal lifetimes.}
% Siblings alias the trunk's KV pages and own private residual pages, yet abort, commit, and the next agent step end those lifetimes at different times.
% Reclamation must free an aborted residual's private pages while any live sibling still references the trunk.
% After commit, the spine stays resident, because the next step uses it as the trunk.

% ---------------------------------------------------------------------------
Achieving this goal, however, presents two key challenges.

\paragraph{Challenge 1: Disaggregated residency.}
Prefill and decode maintain KV on separate instances, making cache residency inherently stage-specific.
A page resident on the prefill instance may still be absent from decode and require connector transfer, while decode-side residency can eliminate that transfer.
The challenge is therefore to efficiently coordinate residency information across the disaggregated prefill and decode instances.

\paragraph{Challenge 2: Prefill-side admission.}
Once the prefill kernel is launched, the uncached residual suffix is computed and decoder-missing KV pages enter the transfer path.
Admission therefore needs to be decided beforehand, using only the residual text and the two residency lengths.
The challenge is to enable reliable prefill-side admission under such limited information with minimal serving-system overhead.

% \paragraph{Challenge 3: Efficiency--quality trade-off.}
% Rejecting a low-scoring residual saves its prefill, its transfer, and its decode, yet that residual may still produce the committed answer.
% Keeping the whole fan-out preserves committed-answer accuracy and pays for every sibling.
% Residuals the pre-execution score leaves uncertain need a short subsequent generation before the system spends a full decode budget.

% ---------------------------------------------------------------------------
% §5 Disaggregated pruning + contributions
% ---------------------------------------------------------------------------
\paragraph{Disaggregated pruning.}
Our key insight is to disaggregate pruning across the prefill and decode stages (\Cref{fig:system}).
\sys admits or rejects residuals before prefill from dual-side KV residency and a lightweight score (\S\ref{sec:app}), refines uncertain residuals with a short decode-side probe (\S\ref{sec:split}), and uses a pruning-aware copy-on-write context tree to share the trunk, reclaim aborted pages, and retain the committed spine (\S\ref{sec:kvstate}).
Existing decode-time pruners operate unchanged on the admitted set.

\paragraph{Results.}
On Qwen3 models, prefill-side admission reduces computed prefill by $23\%$ at matched committed-answer accuracy; on MATH-500, a $128$-token probe restores coverage to $20/32$ at $0.69\times$ the wall-clock of the full fan-out; and spine retention reduces median peak KV usage to $0.72\times$ that of APC (\S\ref{sec:eval}).

\paragraph{Contributions.}
Our contributions are summarized as follows:
\begin{itemize}[leftmargin=*,itemsep=1pt,topsep=2pt]
\item We formulate pruning in disaggregated agentic serving as a per-residual, cross-stage decision problem, and characterize cache regimes by trunk-amplification multiplicity on the prefill instance and on the connector (\S\ref{sec:problem}).
\item We propose \emph{disaggregated pruning} and instantiate it in \sys{} with prefill-side admission and decode-side probing; lightweight KV state management supports branch reclamation and spine retention (\S\ref{sec:design}).
\item We demonstrate that prefill-side admission composes with three existing decode-time pruning methods and reduces computed prefill, and that decode-side probing mitigates the coverage loss of aggressive prefill-side rejection, while preserving committed-answer accuracy (\S\ref{sec:eval}).
\end{itemize}
